% Lösungen Einheit 2 von 3 – Auswahlen und Binomialkoeffizient (zusammenbau v0.1)
\einheitenkopf{Einheit 2 von 3 · Auswahlen und Binomialkoeffizient}

\erg{10}{a) nein – derselbe Becher, gleich in welcher Folge die Sorten gewählt werden \quad b) AB, AC, AD, BC, BD, CD; Anzahl $= 6 = \begin{pmatrix} 4 \\ 2 \end{pmatrix}$ \quad c) $\begin{pmatrix} 40 \\ 6 \end{pmatrix} = 3\,838\,380$ \quad d) gleich, weil zwei wählen sieben liegen lassen heißt; $\begin{pmatrix} 9 \\ 2 \end{pmatrix} = \frac{9 \cdot 8}{2} = 36$ \quad e) Auswahl $\begin{pmatrix} 9 \\ 4 \end{pmatrix} = 126$; Verteilung $4! = 24$ \quad f) zwei Mädchen und ein Junge: $\begin{pmatrix} 5 \\ 2 \end{pmatrix} \cdot 4 = 40$; drei Mädchen: $\begin{pmatrix} 5 \\ 3 \end{pmatrix} = 10$; zusammen $= 50$ \quad g) Lagen der zwei Zusatzzeichen $\begin{pmatrix} 6 \\ 2 \end{pmatrix} = 15$, Zeichenwahl ohne die vier Buchstaben und ohne Wiederholung $46 \cdot 45$; Anzahl $= 15 \cdot 46 \cdot 45 = 31\,050$ \quad h) $\begin{pmatrix} 4 \\ 2 \end{pmatrix} = 6$ wählt die zwei Farben aus vier; 8 zählt die Aufteilungen der neun Rosen auf die zwei Farben – die erste Farbe hat 1 bis 8 Rosen (8 ist nicht die Zahl übriger Rosen) \quad i) Kombination mit Wiederholung: $\begin{pmatrix} 4 + 6 - 1 \\ 6 \end{pmatrix} = \begin{pmatrix} 9 \\ 6 \end{pmatrix} = 84$ (nicht $4^6$)}
\erg{11}{a) Fehler: Ole zählt die Reihenfolge mit; jede Auswahl kommt doppelt vor: $\begin{pmatrix} 8 \\ 2 \end{pmatrix} = 56 : 2 = 28$ \quad b) Wer $k$ Personen auswählt, lässt zugleich die übrigen $n - k$ liegen; jede Auswahl entspricht genau einer liegen gelassenen Gruppe. \quad c) $\begin{pmatrix} 25 \\ 4 \end{pmatrix} = 12\,650$ Auswahlen, eine davon gewinnt: $P = \frac{1}{12\,650}$}
